% !TEX root = ../proyect.tex

\chapter{Dise\~no de la Base de Datos}\label{cap:diseno-bd}

Este cap\'itulo describe el dise\~no de la capa de persistencia del sistema. Se ha optado por MongoDB, una base de datos NoSQL orientada a documentos, accedida a trav\'es de Mongoose como ODM (\textit{Object Document Mapper}). A continuaci\'on se justifica la elecci\'on, se detallan las colecciones dise\~nadas y se describe la estrategia de indexaci\'on adoptada.

\section{Justificaci\'on de MongoDB}\label{sec:justificacion-mongo}

La elecci\'on de MongoDB frente a un sistema relacional tradicional se fundamenta en los siguientes criterios:

\begin{itemize}
    \item \textbf{Flexibilidad del esquema:} Las entidades del dominio, como el horario semanal de un profesional, se modelan de forma natural mediante subdocumentos anidados, evitando m\'ultiples tablas auxiliares y \textit{joins} complejos.
    \item \textbf{Cohesi\'on con el \textit{stack}:} MongoDB se integra nativamente con Node.js y Express, formando parte del \textit{stack} MERN. Los documentos JSON almacenados se transfieren directamente entre base de datos, servidor y cliente sin transformaciones intermedias~\cite{mongodb_docs}.
    \item \textbf{Rendimiento en lecturas:} Las consultas m\'as frecuentes del sistema (disponibilidad de huecos, listado de citas de un d\'ia, fichas de cliente) se resuelven sobre una \'unica colecci\'on con \'indices compuestos, sin necesidad de operaciones de \textit{join}.
    \item \textbf{Escalabilidad horizontal:} Aunque el alcance actual es \textit{single-tenant}, MongoDB facilita el escalado futuro mediante \textit{sharding} y r\'eplicas.
\end{itemize}

No obstante, se mantiene la integridad referencial a nivel de aplicaci\'on mediante Mongoose, que proporciona validaciones de esquema, \textit{hooks} pre/post-guardado y poblaci\'on de referencias (\texttt{populate}).

\section{Modelo de Datos}\label{sec:modelo-datos}

El sistema se compone de siete colecciones principales, representadas en la figura~\ref{fig:modelo-datos}. Cada colecci\'on se describe a continuaci\'on con sus campos, tipos, restricciones y relaciones.

\figura{1.0}{img/modelo_datos}{Modelo de datos del sistema}{fig:modelo-datos}{}

\subsection{Colecci\'on \texttt{users}}\label{subsec:users}

Almacena tanto los clientes que realizan reservas como los administradores del sistema. Se distinguen mediante el campo \texttt{role}. El cuadro~\ref{tab:users} detalla su estructura.

\cuadro{|l|l|l|p{5cm}|}{Estructura de la colecci\'on \texttt{users}}{tab:users}{
\textbf{Campo} & \textbf{Tipo} & \textbf{Requerido} & \textbf{Descripci\'on} \\
\hline
nombre & String & S\'i & Nombre completo (m\'ax. 100 caracteres) \\
email & String & S\'i & Correo electr\'onico, \'unico y en min\'usculas \\
password & String & S\'i & Contrase\~na cifrada con bcrypt (m\'in. 6 caracteres) \\
telefono & String & No & Tel\'efono validado por expresi\'on regular \\
role & String & S\'i & \texttt{cliente} o \texttt{admin} \\
activo & Boolean & S\'i & Permite desactivar cuentas sin eliminarlas \\
passwordChangedAt & Date & No & Marca temporal para invalidar \textit{tokens} anteriores \\
}

La contrase\~na se cifra autom\'aticamente mediante un \textit{hook} \texttt{pre-save} de Mongoose con bcrypt y un factor de coste de 10. El campo \texttt{password} se excluye de las consultas por defecto (\texttt{select: false}) para evitar su exposici\'on accidental. Se define un \'indice sobre el campo \texttt{role} para acelerar las consultas de listado de clientes.

\subsection{Colecci\'on \texttt{professionals}}\label{subsec:professionals}

Representa a los profesionales (peluqueros/barberos) del establecimiento. Cada profesional tiene un horario semanal configurable con soporte para descansos. El cuadro~\ref{tab:professionals} muestra los campos principales.

\cuadro{|l|l|l|p{5cm}|}{Estructura de la colecci\'on \texttt{professionals}}{tab:professionals}{
\textbf{Campo} & \textbf{Tipo} & \textbf{Requerido} & \textbf{Descripci\'on} \\
\hline
nombre & String & S\'i & Nombre del profesional (m\'ax. 100) \\
especialidad & String & S\'i & Especialidad principal (m\'ax. 100) \\
color & String & No & Color hexadecimal para el calendario (por defecto \texttt{\#3B82F6}) \\
activo & Boolean & S\'i & Indica si el profesional est\'a disponible \\
horarioSemanal & Object & S\'i & Subdocumento con la configuraci\'on de cada d\'ia \\
}

El campo \texttt{horarioSemanal} es un subdocumento con siete entradas (claves \texttt{0} a \texttt{6}, donde 0 es domingo), cada una con los campos: \texttt{activo} (Boolean), \texttt{inicio} y \texttt{fin} (String en formato HH:MM), y opcionalmente \texttt{descansoInicio} y \texttt{descansoFin} para la pausa intermedia. Este dise\~no embebido evita una colecci\'on separada de horarios y permite recuperar el horario completo en una sola consulta. El modelo expone m\'etodos de instancia \texttt{getHorarioDia(dia)} y \texttt{trabajaElDia(dia)} que encapsulan la l\'ogica de acceso.

\subsection{Colecci\'on \texttt{services}}\label{subsec:services}

Define el cat\'alogo de servicios ofrecidos por el establecimiento. El cuadro~\ref{tab:services} presenta su estructura.

\cuadro{|l|l|l|p{5cm}|}{Estructura de la colecci\'on \texttt{services}}{tab:services}{
\textbf{Campo} & \textbf{Tipo} & \textbf{Requerido} & \textbf{Descripci\'on} \\
\hline
nombre & String & S\'i & Nombre del servicio (m\'ax. 100) \\
descripcion & String & No & Descripci\'on detallada (m\'ax. 500) \\
duracion & Number & S\'i & Duraci\'on en minutos (m\'in. 15) \\
precio & Number & S\'i & Precio en euros (m\'in. 0) \\
categoria & String & No & Categor\'ia para agrupaci\'on (m\'ax. 50) \\
profesionalesCapaces & [ObjectId] & No & Referencias a profesionales que ofrecen el servicio \\
activo & Boolean & S\'i & Visibilidad en el cat\'alogo p\'ublico \\
orden & Number & No & Posici\'on para ordenaci\'on en la interfaz \\
}

El campo \texttt{profesionalesCapaces} establece una relaci\'on muchos-a-muchos entre servicios y profesionales, almacenada como un \textit{array} de referencias. Esta decisi\'on permite que el motor de disponibilidad filtre r\'apidamente qu\'e profesionales pueden atender un servicio concreto. El modelo incluye propiedades virtuales \texttt{duracionFormateada} y \texttt{precioFormateado} que generan cadenas legibles (por ejemplo, ``1h 30min'' o ``25,50 \euro{}'').

\subsection{Colecci\'on \texttt{appointments}}\label{subsec:appointments}

Es la colecci\'on central del sistema, que registra cada cita reservada. El cuadro~\ref{tab:appointments} detalla sus campos.

\cuadro{|l|l|l|p{4.5cm}|}{Estructura de la colecci\'on \texttt{appointments}}{tab:appointments}{
\textbf{Campo} & \textbf{Tipo} & \textbf{Requerido} & \textbf{Descripci\'on} \\
\hline
cliente & ObjectId & S\'i & Referencia al usuario que reserva \\
profesional & ObjectId & S\'i & Referencia al profesional asignado \\
servicio & ObjectId & S\'i & Referencia al servicio solicitado \\
fechaHoraInicio & Date & S\'i & Inicio de la cita (UTC) \\
fechaHoraFin & Date & S\'i & Fin real de la cita (UTC) \\
fechaHoraFinOperativa & Date & No & Fin operativo redondeado a la rejilla de \textit{slots} \\
duracionOperativaMinutos & Number & No & Duraci\'on operativa para bloqueo de agenda (m\'in. 15) \\
estado & String & S\'i & \texttt{confirmada}, \texttt{completada}, \texttt{cancelada} o \texttt{no\_presentado} \\
precioFinal & Number & No & Instant\'anea del precio en el momento de la reserva \\
notasCliente & String & No & Observaciones del cliente (m\'ax. 500) \\
notasInternas & String & No & Notas privadas del administrador (m\'ax. 1000) \\
forzadaPorAdmin & Boolean & No & Indica si se permiti\'o \textit{overbooking} \\
motivoCancelacion & String & No & Raz\'on de la cancelaci\'on (m\'ax. 500) \\
canceladaPor & String & No & \texttt{cliente} o \texttt{admin} \\
}

Se distingue entre \texttt{fechaHoraFin} (duraci\'on real del servicio) y \texttt{fechaHoraFinOperativa} (duraci\'on redondeada a la rejilla de 15 o 30 minutos configurada en los ajustes del negocio). Esta distinci\'on permite que un servicio de 20 minutos bloquee un \textit{slot} de 30 minutos en la agenda, evitando fragmentaci\'on. El modelo incluye propiedades virtuales \texttt{esPasada} y \texttt{esHoy} que facilitan el filtrado temporal, y un m\'etodo \texttt{puedeCancelar(horasMinimas)} que verifica si la cita admite cancelaci\'on seg\'un la pol\'itica del negocio.

Esta colecci\'on cuenta con m\'ultiples \'indices compuestos, detallados en la secci\'on~\ref{sec:indices}.

\subsection{Colecci\'on \texttt{blockers}}\label{subsec:blockers}

Registra los periodos de tiempo en los que no se pueden agendar citas, ya sea para un profesional concreto o para todo el negocio. El cuadro~\ref{tab:blockers} muestra su estructura.

\cuadro{|l|l|l|p{4.5cm}|}{Estructura de la colecci\'on \texttt{blockers}}{tab:blockers}{
\textbf{Campo} & \textbf{Tipo} & \textbf{Requerido} & \textbf{Descripci\'on} \\
\hline
profesional & ObjectId & No & Profesional afectado (\texttt{null} = bloqueo global) \\
titulo & String & S\'i & T\'itulo descriptivo (m\'ax. 100) \\
descripcion & String & No & Detalle del bloqueo (m\'ax. 500) \\
fechaHoraInicio & Date & S\'i & Inicio del bloqueo (UTC) \\
fechaHoraFin & Date & S\'i & Fin del bloqueo (debe ser posterior al inicio) \\
tipo & String & S\'i & \texttt{vacaciones}, \texttt{festivo}, \texttt{personal}, \texttt{mantenimiento} u \texttt{otro} \\
esRecurrente & Boolean & No & Reservado para funcionalidad futura \\
creadoPor & ObjectId & No & Referencia al administrador que lo cre\'o \\
}

Cuando el campo \texttt{profesional} es \texttt{null}, el bloqueo se aplica a todo el negocio (por ejemplo, un d\'ia festivo). El modelo expone m\'etodos est\'aticos \texttt{hayBloqueo(profesionalId, inicio, fin)} para verificaci\'on r\'apida de solapamientos y \texttt{getBloqueos(profesionalId, inicio, fin)} para obtener todos los bloqueos aplicables a un rango, incluyendo los globales.

\subsection{Colecci\'on \texttt{technicalnotes}}\label{subsec:technicalnotes}

Implementa la funcionalidad de CRM del sistema, permitiendo a los profesionales registrar informaci\'on t\'ecnica sobre cada cliente. El cuadro~\ref{tab:technicalnotes} detalla sus campos.

\cuadro{|l|l|l|p{4.5cm}|}{Estructura de la colecci\'on \texttt{technicalnotes}}{tab:technicalnotes}{
\textbf{Campo} & \textbf{Tipo} & \textbf{Requerido} & \textbf{Descripci\'on} \\
\hline
cliente & ObjectId & S\'i & Referencia al cliente \\
cita & ObjectId & No & Referencia a la cita asociada (opcional) \\
creadaPor & ObjectId & No & Referencia al profesional autor \\
titulo & String & No & T\'itulo de la nota (m\'ax. 100) \\
contenido & String & S\'i & Contenido de la nota (m\'ax. 2000) \\
categoria & String & No & \texttt{color}, \texttt{tratamiento}, \texttt{alergia}, \texttt{preferencia} u \texttt{otro} \\
importante & Boolean & No & Marca para notas cr\'iticas (alergias, etc.) \\
}

Las categor\'ias permiten clasificar las notas seg\'un su naturaleza: f\'ormulas de \textit{color} aplicadas, \textit{tratamientos} realizados, \textit{alergias} detectadas o \textit{preferencias} del cliente. La marca \texttt{importante} permite destacar informaci\'on cr\'itica que debe consultarse antes de cada servicio.

\subsection{Colecci\'on \texttt{settings}}\label{subsec:settings}

Almacena la configuraci\'on global del negocio en un \'unico documento con \texttt{\_id} fijo igual a \texttt{global} (patr\'on \textit{singleton}). El cuadro~\ref{tab:settings} presenta los campos configurables.

\cuadro{|l|l|p{6cm}|}{Estructura de la colecci\'on \texttt{settings}}{tab:settings}{
\textbf{Campo} & \textbf{Tipo} & \textbf{Descripci\'on} \\
\hline
nombreNegocio & String & Nombre del establecimiento \\
telefono & String & Tel\'efono de contacto \\
email & String & Correo electr\'onico del negocio \\
direccion & String & Direcci\'on f\'isica \\
horasMinimasCancelacion & Number & Horas de antelaci\'on m\'inima para cancelar (por defecto 24) \\
diasMaximosReserva & Number & D\'ias m\'aximos de antelaci\'on para reservar (1--365) \\
duracionSlot & Number & Rejilla de \textit{slots}: 15 o 30 minutos \\
zonaHoraria & String & Zona horaria IANA (por defecto \texttt{Europe/Madrid}) \\
mensajeBienvenida & String & Mensaje personalizado en el portal p\'ublico \\
politicaCancelacion & String & Texto de la pol\'itica de cancelaci\'on \\
}

Este documento se carga en memoria al arrancar la aplicaci\'on y se almacena en cach\'e con \texttt{Object.freeze()} para garantizar inmutabilidad. Los m\'etodos est\'aticos \texttt{getGlobal()} y \texttt{updateGlobal()} gestionan el acceso y la invalidaci\'on de la cach\'e respectivamente, evitando consultas repetidas a la base de datos en cada petici\'on.

\section{Estrategia de Indexaci\'on}\label{sec:indices}

Se han definido \'indices para optimizar las consultas m\'as frecuentes del sistema. El cuadro~\ref{tab:indices} resume los \'indices principales.

\cuadro{|l|l|p{5cm}|}{\'Indices principales definidos en las colecciones}{tab:indices}{
\textbf{Colecci\'on} & \textbf{\'Indice} & \textbf{Prop\'osito} \\
\hline
users & \{role: 1\} & Filtrar clientes vs. administradores \\
appointments & \{profesional: 1, fechaHoraInicio: 1\} & Agenda diaria de un profesional \\
appointments & \{cliente: 1, fechaHoraInicio: -1\} & Historial de citas de un cliente \\
appointments & \{fechaHoraInicio: 1, fechaHoraFin: 1\} & Consultas de disponibilidad por rango \\
appointments & \{estado: 1\} & Filtrado por estado de la cita \\
professionals & \{activo: 1\} & Listado de profesionales activos \\
services & \{activo: 1, orden: 1\} & Cat\'alogo ordenado de servicios activos \\
services & \{profesionalesCapaces: 1\} & B\'usqueda de servicios por profesional \\
blockers & \{profesional: 1, fechaHoraInicio: 1, fechaHoraFin: 1\} & Detecci\'on de bloqueos por profesional y rango \\
technicalnotes & \{cliente: 1, createdAt: -1\} & Historial de notas de un cliente \\
technicalnotes & \{cliente: 1, importante: 1\} & Notas cr\'iticas de un cliente \\
}

Los \'indices sobre \texttt{appointments} son especialmente cr\'iticos, ya que el motor de disponibilidad realiza consultas de rango temporal para cada \textit{slot} calculado. El \'indice compuesto \texttt{\{profesional, fechaHoraInicio\}} permite resolver la consulta de ``citas de un profesional en un d\'ia concreto'' mediante un recorrido secuencial del \'indice (\textit{index scan}) sin acceder a los documentos.

\section{Relaciones entre Colecciones}\label{sec:relaciones}

Aunque MongoDB no soporta \textit{joins} nativos como los sistemas relacionales, las relaciones entre colecciones se implementan mediante referencias (\texttt{ObjectId}) y se resuelven en tiempo de consulta con el m\'etodo \texttt{populate()} de Mongoose~\cite{mongoose_docs}. Las relaciones principales son:

\begin{itemize}
    \item \textbf{Appointment $\rightarrow$ User:} Cada cita referencia al cliente que la reserv\'o.
    \item \textbf{Appointment $\rightarrow$ Professional:} Cada cita referencia al profesional asignado.
    \item \textbf{Appointment $\rightarrow$ Service:} Cada cita referencia al servicio solicitado.
    \item \textbf{Service $\rightarrow$ Professional:} Cada servicio contiene un \textit{array} de referencias a los profesionales capacitados para realizarlo (relaci\'on muchos-a-muchos).
    \item \textbf{Blocker $\rightarrow$ Professional:} Cada bloqueo puede asociarse a un profesional o ser global.
    \item \textbf{TechnicalNote $\rightarrow$ User:} Cada nota t\'ecnica referencia al cliente sobre el que se escribe.
    \item \textbf{TechnicalNote $\rightarrow$ Appointment:} Opcionalmente, una nota puede vincularse a una cita concreta.
\end{itemize}

La integridad referencial se garantiza a nivel de aplicaci\'on: antes de crear una cita, el controlador verifica que el cliente, el profesional y el servicio existen y est\'an activos. Mongoose valida los tipos \texttt{ObjectId} en el esquema, rechazando documentos con referencias mal formadas.

